% =====================================================================
%  DE DOA TOI TINH HOP LE.
%
%  Muc nay ton tai de KHAI TRUOC nhung gi phan bien se hoi. Luat cua nha:
%  ``chua mo hinh'' ma khong kem SO chi la mot loi xin loi. Nen moi han che o
%  day phai kem HOAC mot con so HOAC mot phat bieu kiem duoc.
%
%  ⚠ KHONG GO SO. Xem paper/check-no-typed-numbers.py
% =====================================================================

\section{Threats to Validity}
\label{sec:threats}

\textbf{No constrained hardware.} Every measurement here runs on a desktop over loopback. This
bounds what may be claimed, and the claims are bounded accordingly: exchange counts are a
property of the transport, and no statement is made about latency, energy, or memory. Those
have been measured on real hardware by others, for TLS and DTLS over
6TiSCH~\cite{claeys2021}, for an embedded EDHOC module~\cite{fraile2021}, and for
post-quantum EDHOC~\cite{fraile2025}; this letter cites rather than repeats them.
The one quantity that would change on real hardware, retransmission behaviour under loss, is
addressed only qualitatively.

\textbf{One implementation excluded, not reported as a failure.} mbedTLS cannot load the
largest certificate at all: \texttt{mbedtls\_x509\_crt\_parse\_file} rejects the key, which
exceeds the library's configured maximum integer size. Its cells for that credential would
otherwise read as failures at every MTU, including the widest, and would look exactly like a
result about the link. They are a limit of the library, reached before the link is touched.
The per-pair control described in Section~\ref{sec:setup} exists for this case, and this is
the pair it removed.

\textbf{Findings are about implementations, not about DTLS.} The three libraries are bounded
by different quantities at very different points, so nothing here supports a statement about
the protocol. Where DTLS~1.3 itself is characterised, in particular its constant round-trip
count, the source is the specification \cite{rfc9147} and not these measurements. No DTLS~1.3
implementation was available to us; \texttt{s\_server} in the OpenSSL version used still
offers no DTLS~1.3 option, so the flight behaviour was exercised on DTLS~1.2, where the
mechanism is the same and only the flight count differs.

\textbf{Loss is assumed independent.} The block-size sweep treats frame losses as independent,
which is optimistic for a low-power radio: bursty loss favours larger blocks, and would move
the optimum. The conclusion drawn from that sweep is deliberately weaker than the curve
allows, namely that no permitted block size brings the exchange count near the DTLS flight
count, and this holds across the whole loss range plotted, including the lossless case.

\textbf{One estimated parameter.} The MAC and link-layer security overhead subtracted from the
frame is an estimate rather than a measurement. It affects the frame payload and so the
fragment counts. It cannot affect the threshold argument: the smallest standardised KEM
encapsulation key is $\numKemFrameRatio\times$ the frame payload, and no plausible correction
of a few bytes changes that.

\textbf{Scope is EDHOC over CoAP.} The lock-step behaviour follows from CoAP block-wise
transfer, so it applies to the profile of \cite{rfc9668} and not to EDHOC carried over some
other transport. The per-message decomposition reproduces the published total to within
\numDecompError\%, which is adequate for counting exchanges but not for quoting absolute
sizes; absolute sizes are taken from the published comparison instead.
